%! Multiplying \& Dividing Rational Expressions — Unit 5, Lesson 5.2 — Student Packet Cover Sheet
\documentclass[10pt]{article}
\usepackage{algebra2-article}
\usepackage{algebra2-boxes}
\usepackage{ltablex}
\keepXColumns

\begin{document}

% Full-bleed forest banner
\begin{tikzpicture}[remember picture, overlay]
  \fill[forest] (current page.north west)
    rectangle ([yshift=-0.9in]current page.north east);
\end{tikzpicture}
\vspace{-0.8in}

\begin{center}
  {\color{white}\textbf{\LARGE Algebra 2: Shepherd}}\\[4pt]
  {\color{forestmid}\large\textbf{Unit 5 \quad Rational Functions}}\\[6pt]
  {\color{white}\textbf{\large Lesson 5.2 \quad Multiplying \& Dividing Rational Expressions}}
\end{center}

\vspace{0.22in}
\namedateperiod
\vspace{0.18in}

\begin{learningtargetbox}
\small
\begin{enumerate}[leftmargin=1.5em, itemsep=5pt]
  \item \ldots{} \textbf{multiply} rational expressions by \textbf{factoring first} and dividing out common factors across the whole product --- never by multiplying out.
  \item \ldots{} \textbf{divide} rational expressions by multiplying by the \textbf{reciprocal} of the divisor (keep--change--flip), flipping \emph{before} I cancel anything.
  \item \ldots{} list \textbf{every} restriction from all \textbf{three} sources --- each original denominator \emph{and} the \textbf{numerator of the divisor} --- and say where each one came from.
\end{enumerate}
\end{learningtargetbox}

\vspace{0.14in}

\begin{tocbox}
\small
\renewcommand{\arraystretch}{1.5}
\begin{tabularx}{\linewidth}{@{} c l X r @{}}
\rowcolor{forestbg}
\textbf{\#} & \textbf{Component} & \textbf{Description} & \textbf{Score} \\
\midrule
1 & Warm-Up        & Spiral review: factor-and-cancel with numbers, factoring completely, and yesterday's simplify-and-restrict & \blank{1.2cm} \\
2 & Guided Notes   & Products, keep--change--flip, and the three sources of a restriction & \blank{1.2cm} \\
3 & Group Activity & \emph{Factor First, Flip Second, Restrict Always} --- a feature table, an error analysis, and restriction detective, by tier & \blank{1.2cm} \\
4 & Exit Ticket    & Quick check: a quotient with all three sources, plus an SOL-style item & \blank{1.2cm} \\
5 & Homework       & Monomial and binomial practice, a rectangle model, and a design-your-own extension & \blank{1.2cm} \\
\midrule
  & \multicolumn{2}{r}{\textbf{Total}} & \blank{1.2cm} \\
\end{tabularx}
\end{tocbox}

\vspace{0.10in}

\begin{remindbox}
\small To \textbf{multiply} rational expressions, factor everything, divide out any numerator factor
against any denominator factor across the \emph{whole} product, and multiply what is left --- never
expand first, because an expanded product hides its own factors. To \textbf{divide}, use
\textbf{keep--change--flip}: keep the first fraction, change $\div$ to $\times$, and flip the divisor,
because $\frac cd\cdot\frac dc=1$. Flip \emph{before} you cancel --- the $\div$ symbol is not a
fraction bar, and cancelling across it gives the wrong answer. The new idea today is where the
\textbf{restrictions} come from: a quotient $\frac AB\div\frac CD$ needs $B\neq0$ \emph{and} $D\neq0$
(both fractions have to exist) \emph{and} $C\neq0$, because the divisor equals zero exactly when its
\textbf{numerator} does --- and you can never divide by zero. Source 2 goes invisible after the flip;
source 3 is invisible before it. \textbf{Flipping trades one blind spot for another}, so read the
restrictions off the problem as written, then add the divisor's numerator. And as always, an answer
with no restriction list is half an answer --- even when the answer is a polynomial.
\end{remindbox}

\end{document}
